%==============================================================================
\section{Chương 1. Cơ sở Gradient Descent}
\label{sec:gd-part1}
%==============================================================================

\subsection{1.1. Giới thiệu}

\begin{ynghia}[title={Bề mặt hàm số và cực tiểu}]
Hình minh họa dưới đây thể hiện đồ thị hàm một biến và vị trí cực tiểu:

\mlcbimg{01_gradient_descent_part1/img02.png}

Điểm màu xanh lục là điểm \textbf{local minimum} (cực tiểu), và cũng là điểm làm cho hàm số đạt giá trị nhỏ nhất. Từ đây trở đi, ta dùng \emph{local minimum} để thay cho điểm cực tiểu, \emph{global minimum} để thay cho điểm mà tại đó hàm số đạt giá trị nhỏ nhất. \textbf{Global minimum} là một trường hợp đặc biệt của local minimum.
\end{ynghia}

\begin{dinhnghia}[title={Local minimum và đạo hàm}]
Giả sử chúng ta đang quan tâm đến một hàm số một biến có đạo hàm mọi nơi. Nhắc lại vài điều đã quá quen thuộc:

Điểm local minimum $x^*$ của hàm số là điểm có đạo hàm $f'(x^*)$ bằng $0$. Hơn thế nữa, trong lân cận của nó, đạo hàm của các điểm phía bên trái $x^*$ là không dương, đạo hàm của các điểm phía bên phải $x^*$ là không âm.

Đường tiếp tuyến với đồ thị hàm số đó tại một điểm bất kỳ có hệ số góc chính bằng đạo hàm của hàm số tại điểm đó.
\end{dinhnghia}

\begin{tinhchat}[title={Dấu đạo hàm quanh cực tiểu}]
Trong hình phía trên, các điểm bên trái của điểm local minimum màu xanh lục có đạo hàm âm, các điểm bên phải có đạo hàm dương. Và đối với hàm số này, càng xa về phía trái của điểm local minimum thì đạo hàm càng âm, càng xa về phía phải thì đạo hàm càng dương.
\end{tinhchat}

\subsubsection{Gradient Descent}

\begin{dongco}[title={Tại sao cần tối ưu trong Machine Learning?}]
Trong Machine Learning nói riêng và Toán Tối Ưu nói chung, chúng ta thường xuyên phải tìm giá trị nhỏ nhất (hoặc đôi khi là lớn nhất) của một hàm số nào đó. Ví dụ như các hàm mất mát trong hai bài Linear Regression và K-means Clustering. Nhìn chung, việc tìm global minimum của các hàm mất mát trong Machine Learning là rất phức tạp, thậm chí là bất khả thi. Thay vào đó, người ta thường cố gắng tìm các điểm local minimum, và ở một mức độ nào đó, coi đó là nghiệm cần tìm của bài toán.
\end{dongco}

\begin{phuongphap}[title={Từ phương trình đạo hàm bằng 0 đến lặp}]
Các điểm local minimum là nghiệm của phương trình đạo hàm bằng $0$. Nếu bằng một cách nào đó có thể tìm được toàn bộ (hữu hạn) các điểm cực tiểu, ta chỉ cần thay từng điểm local minimum đó vào hàm số rồi tìm điểm làm cho hàm có giá trị nhỏ nhất (bước so sánh giá trị hàm tại các cực tiểu). Tuy nhiên, trong hầu hết các trường hợp, việc giải phương trình đạo hàm bằng $0$ là bất khả thi. Nguyên nhân có thể đến từ sự phức tạp của dạng của đạo hàm, từ việc các điểm dữ liệu có số chiều lớn, hoặc từ việc có quá nhiều điểm dữ liệu.

Hướng tiếp cận phổ biến nhất là xuất phát từ một điểm mà chúng ta coi là gần với nghiệm của bài toán, sau đó dùng một phép toán lặp để tiến dần đến điểm cần tìm, tức đến khi đạo hàm gần với $0$. \textbf{Gradient Descent} (viết gọn là GD) và các biến thể của nó là một trong những phương pháp được dùng nhiều nhất.
\end{phuongphap}

\begin{phuongphap}[title={Cấu trúc chương}]
Nội dung GD được chia hai chương: \textbf{Chương~1} (mục này) --- ý tưởng, cập nhật một/nhiều biến, Linear Regression và kiểm tra gradient; \textbf{Chương~2} --- momentum, NAG, Batch/SGD/mini-batch, điều kiện dừng, Newton. Bài toán \emph{large-scale} (số chiều và số mẫu lớn) là động lực cho các biến thể ở Chương~2.
\end{phuongphap}

\begin{phuongphap}[title={Thuật toán Gradient Descent (khung chung)}]
\textbf{Đầu vào:} hàm khả vi $f(\boldsymbol{\theta})$ (hoặc $\mathcal{L}$), learning rate $\eta > 0$, điểm khởi tạo $\boldsymbol{\theta}_0$, tiêu chí dừng.\\
\textbf{Lặp} cho $t = 0, 1, 2, \ldots$ đến khi dừng:
\begin{enumerate}
  \item Tính $\mathbf{g}_t = \nabla_{\boldsymbol{\theta}} f(\boldsymbol{\theta}_t)$.
  \item Cập nhật $\boldsymbol{\theta}_{t+1} = \boldsymbol{\theta}_t - \eta\, \mathbf{g}_t$.
  \item Nếu $\|\mathbf{g}_t\|$ nhỏ hơn ngưỡng (hoặc đạt số vòng lặp tối đa) thì dừng.
\end{enumerate}
\textbf{Đầu ra:} $\boldsymbol{\theta}_t$ xấp xỉ cực tiểu (thường là local minimum).
\end{phuongphap}

%------------------------------------------------------------------------------
\subsection{1.2. Gradient Descent cho hàm một biến}
%------------------------------------------------------------------------------

\begin{ynghia}[title={Quan sát từ đồ thị}]
Quay trở lại hình vẽ ban đầu và các quan sát đã nêu. Giả sử $x_t$ là điểm ta tìm được sau vòng lặp thứ $t$. Ta cần tìm một thuật toán để đưa $x_t$ về càng gần $x^*$ càng tốt.

Trong hình đầu tiên, chúng ta lại có thêm hai quan sát nữa:

Nếu đạo hàm của hàm số tại $x_t$: $f'(x_t) > 0$ thì $x_t$ nằm về bên phải so với $x^*$ (và ngược lại). Để điểm tiếp theo $x_{t+1}$ gần với $x^*$ hơn, chúng ta cần di chuyển $x_t$ về phía bên trái, tức về phía âm. Nói cách khác, chúng ta cần di chuyển ngược dấu với đạo hàm:
\[
x_{t+1} = x_t + \Delta
\]
Trong đó $\Delta$ là một đại lượng ngược dấu với đạo hàm $f'(x_t)$.

$x_t$ càng xa $x^*$ về phía bên phải thì $f'(x_t)$ càng lớn hơn $0$ (và ngược lại). Vậy, lượng di chuyển $\Delta$, một cách trực quan nhất, là tỉ lệ thuận với $-f'(x_t)$.
\end{ynghia}

\begin{congthuc}[title={Cập nhật Gradient Descent một biến}]
Hai nhận xét phía trên cho ta công thức cập nhật:
\[
x_{t+1} = x_t - \eta\, f'(x_t)
\]
Trong đó $\eta$ (đọc là \emph{eta}) là \textbf{learning rate} (bước học), số dương. Dấu \textbf{trừ} vì ta đi \textbf{ngược} hướng đạo hàm --- gradient chỉ hướng \textbf{tăng} nhanh nhất của $f$, còn ta cần \textbf{giảm} $f$ (Gradient \emph{Descent} = ``đi xuống'').
\end{congthuc}

\begin{ynghia}[title={Điều kiện giảm hàm (một biến, $\eta$ nhỏ)}]
Với $f$ khả vi, nếu $f'(x_t) \neq 0$ và $|\eta f'(x_t)|$ đủ nhỏ thì $f(x_{t+1}) < f(x_t)$ (xấp xỉ bậc nhất). Đây là cơ sở lý thuyết cho việc chọn $\eta$ không quá lớn.
\end{ynghia}

\subsubsection{Ví dụ đơn giản với Python}

\begin{vidu}[title={Hàm $f(x) = x^2 + 5\sin(x)$}]
Xét hàm số $f(x) = x^2 + 5\sin(x)$ với đạo hàm $f'(x) = 2x + 5\cos(x)$ (hàm này được chọn vì nó không dễ tìm nghiệm của đạo hàm bằng $0$ như hàm phía trên). Giả sử bắt đầu từ một điểm $x_0$ nào đó, tại vòng lặp thứ $t$, chúng ta sẽ cập nhật như sau:
\[
x_{t+1} = x_t - \eta\bigl(2x_t + 5\cos(x_t)\bigr)
\]

Đầu tiên, ta khai báo vài thư viện quen thuộc:
\begin{lstlisting}[language=Python]
# To support both python 2 and python 3
from __future__ import division, print_function, unicode_literals
import math
import numpy as np
import matplotlib.pyplot as plt
\end{lstlisting}

Tiếp theo, ta định nghĩa các hàm số:
\begin{enumerate}
 \item \texttt{grad} để tính đạo hàm;
 \item \texttt{cost} để tính giá trị của hàm số. Hàm này không sử dụng trong thuật toán nhưng thường được dùng để kiểm tra việc tính đạo hàm đúng không hoặc để xem giá trị của hàm số có giảm theo mỗi vòng lặp hay không;
 \item \texttt{myGD1} là phần chính thực hiện thuật toán Gradient Descent nêu phía trên. Đầu vào của hàm số này là learning rate và điểm bắt đầu. Thuật toán dừng lại khi đạo hàm có độ lớn đủ nhỏ.
\end{enumerate}

\begin{lstlisting}[language=Python]
def grad(x):
 return 2*x+ 5*np.cos(x)

def cost(x):
 return x**2 + 5*np.sin(x)

def myGD1(eta, x0):
 x = [x0]
 for it in range(100):
 x_new = x[-1] - eta*grad(x[-1])
 if abs(grad(x_new)) < 1e-3:
 break
 x.append(x_new)
 return (x, it)
\end{lstlisting}
\end{vidu}

\paragraph{Điểm khởi tạo khác nhau.}

\begin{vidu}[title={Thử $x_0 = -5$ và $x_0 = 5$}]
Sau khi có các hàm cần thiết, thử nghiệm với nghiệm với các điểm khởi tạo khác nhau là $x_0 = -5$ và $x_0 = 5$:
\begin{lstlisting}[language=Python]
(x1, it1) = myGD1(.1, -5)
(x2, it2) = myGD1(.1, 5)
print('Solution x1 = %f, cost = %f, obtained after %d iterations'%(x1[-1], cost(x1[-1]), it1))
print('Solution x2 = %f, cost = %f, obtained after %d iterations'%(x2[-1], cost(x2[-1]), it2))
\end{lstlisting}

Kết quả in ra:
\begin{lstlisting}
Solution x1 = -1.110667, cost = -3.246394, obtained after 11 iterations
Solution x2 = -1.110341, cost = -3.246394, obtained after 29 iterations
\end{lstlisting}

Vậy là với các điểm ban đầu khác nhau, thuật toán của chúng ta tìm được nghiệm gần giống nhau, mặc dù với tốc độ hội tụ khác nhau. Dưới đây là hình ảnh minh họa thuật toán GD cho bài toán này .

\mlcbimg{01_gradient_descent_part1/img03.png}
\mlcbimg{01_gradient_descent_part1/img04.png}
\end{vidu}

\begin{ynghia}[title={So sánh hai điểm khởi tạo}]
Từ hình minh họa trên ta thấy rằng ở hình bên trái, tương ứng với $x_0 = -5$, nghiệm hội tụ nhanh hơn, vì điểm ban đầu $x_0$ gần với nghiệm $x^* \approx -1$ hơn. Hơn nữa, với $x_0 = 5$ ở hình bên phải, đường đi của nghiệm có chứa một khu vực có đạo hàm khá nhỏ gần điểm có hoành độ bằng $2$. Điều này khiến cho thuật toán la cà ở đây khá lâu. Khi vượt qua được điểm này thì mọi việc diễn ra rất tốt đẹp.
\end{ynghia}

\paragraph{Learning rate khác nhau.}

\begin{vidu}[title={Cùng $x_0 = -5$, $\eta$ khác nhau}]
Tốc độ hội tụ của GD không những phụ thuộc vào điểm khởi tạo ban đầu mà còn phụ thuộc vào learning rate. Dưới đây là một ví dụ với cùng điểm khởi tạo $x_0 = -5$ nhưng learning rate khác nhau:

\mlcbimg{01_gradient_descent_part1/img05.png}
\mlcbimg{01_gradient_descent_part1/img06.png}
\end{vidu}

\begin{tinhchat}[title={Hai trường hợp learning rate}]
Ta quan sát thấy hai điều:
\begin{enumerate}
 \item Với learning rate nhỏ $\eta = 0.01$, tốc độ hội tụ rất chậm. Trong ví dụ này giới hạn tối đa $100$ vòng lặp nên thuật toán dừng lại trước khi tới đích, mặc dù đã rất gần. Trong thực tế, khi việc tính toán trở nên phức tạp, learning rate quá thấp sẽ ảnh hưởng tới tốc độ của thuật toán rất nhiều, thậm chí không bao giờ tới được đích.
 \item Với learning rate lớn $\eta = 0.5$, thuật toán tiến rất nhanh tới gần đích sau vài vòng lặp. Tuy nhiên, thuật toán không hội tụ được vì bước nhảy quá lớn, khiến nó cứ quẩn quanh ở đích.
\end{enumerate}
\end{tinhchat}

\begin{luuy}[title={Chọn learning rate}]
Việc lựa chọn learning rate rất quan trọng trong các bài toán thực tế. Việc lựa chọn giá trị này phụ thuộc nhiều vào từng bài toán và phải làm một vài thí nghiệm để chọn ra giá trị tốt nhất. Ngoài ra, tùy vào một số bài toán, GD có thể làm việc hiệu quả hơn bằng cách chọn ra learning rate phù hợp hoặc chọn learning rate khác nhau ở mỗi vòng lặp. Phần sau trình bày vấn đề này ở phần 2.
\end{luuy}

%------------------------------------------------------------------------------
\subsection{1.3. Gradient Descent cho hàm nhiều biến}
%------------------------------------------------------------------------------

\begin{dinhnghia}[title={Cập nhật với vector tham số $\boldsymbol{\theta}$}]
Giả sử ta cần tìm global minimum cho hàm $f(\boldsymbol{\theta})$ trong đó $\boldsymbol{\theta}$ (\emph{theta}) là một vector, thường được dùng để ký hiệu tập hợp các tham số của một mô hình cần tối ưu (trong Linear Regression thì các tham số chính là hệ số $\mathbf{w}$). Đạo hàm của hàm số đó tại một điểm $\boldsymbol{\theta}$ bất kỳ được ký hiệu là $\nabla_{\boldsymbol{\theta}} f(\boldsymbol{\theta})$ (hình tam giác ngược đọc là \emph{nabla}). Tương tự như hàm một biến, thuật toán GD cho hàm nhiều biến cũng bắt đầu bằng một điểm dự đoán $\boldsymbol{\theta}_0$, sau đó, ở vòng lặp thứ $t$, quy tắc cập nhật là:
\end{dinhnghia}

\begin{congthuc}[title={Quy tắc cập nhật nhiều biến}]
\[
\boldsymbol{\theta}_{t+1} = \boldsymbol{\theta}_t - \eta\, \nabla_{\boldsymbol{\theta}} f(\boldsymbol{\theta}_t)
\]
Hoặc viết dưới dạng đơn giản hơn: $\boldsymbol{\theta} \leftarrow \boldsymbol{\theta} - \eta\, \nabla_{\boldsymbol{\theta}} f(\boldsymbol{\theta})$.

\textbf{Quy tắc cần nhớ:} luôn luôn đi ngược hướng với đạo hàm.

Việc tính toán đạo hàm của các hàm nhiều biến là một kỹ năng cần thiết.
\end{congthuc}

\subsubsection{Ứng dụng: Linear Regression}

\begin{ynghia}[title={Tối ưu hàm mất mát Linear Regression}]
Áp dụng GD cho bài toán Linear Regression: tối ưu hàm mất mát thay vì (hoặc để đối chiếu với) nghiệm đóng $\mathbf{w} = (\bar{\mathbf{X}}^{\mathsf{T}}\bar{\mathbf{X}})^{-1}\bar{\mathbf{X}}^{\mathsf{T}}\mathbf{y}$.
\end{ynghia}

\begin{dongco}[title={GD so với nghiệm công thức}]
Với MSE và ma trận đầy đủ hạng, nghiệp vụ tối thiểu là duy nhất (hàm lồi). GD và công thức ma trận \textbf{cùng hướng tới cùng $\mathbf{w}^*$} nếu $\eta$ phù hợp. GD vẫn quan trọng vì: (i) không cần lưu nghịch đảo ma trận lớn; (ii) mở rộng được khi thêm regularization hoặc mô hình phi tuyến; (iii) là mẫu cho SGD khi $N$ rất lớn.
\end{dongco}

\begin{congthuc}[title={Hàm mất mát và gradient}]
Hàm mất mát của Linear Regression là:
\[
\mathcal{L}(\mathbf{w}) = \frac{1}{2N}\,\|\mathbf{y} - \bar{\mathbf{X}}\mathbf{w}\|_2^2
\]

Đạo hàm của hàm mất mát là:
\[
\nabla_{\mathbf{w}}\mathcal{L}(\mathbf{w}) = \frac{1}{N}\,\bar{\mathbf{X}}^{\mathsf{T}}\bigl(\bar{\mathbf{X}}\mathbf{w} - \mathbf{y}\bigr) \qquad (1)
\]
\end{congthuc}

\begin{luuy}[title={Hàm mất mát có chia cho $N$}]
Hàm này có khác một chút so với hàm trong bài Linear Regression (không chia $N$). Mẫu số có thêm $N$ là số lượng dữ liệu trong training set. Việc lấy trung bình cộng của lỗi này nhằm giúp tránh trường hợp hàm mất mát và đạo hàm có giá trị là một số rất lớn, ảnh hưởng tới độ chính xác của các phép toán khi thực hiện trên máy tính. Về mặt toán học, nghiệm của hai bài toán là như nhau.
\end{luuy}

\subsubsection{Ví dụ trên Python và lưu ý khi lập trình}

\begin{vidu}[title={Load thư viện và tạo dữ liệu}]
\textbf{Load thư viện:}
\begin{lstlisting}[language=Python]
# To support both python 2 and python 3
from __future__ import division, print_function, unicode_literals
import numpy as np
import matplotlib
import matplotlib.pyplot as plt
np.random.seed(2)
\end{lstlisting}

\textbf{Tiếp theo,} chúng ta tạo $1000$ điểm dữ liệu được chọn gần với đường thẳng $y = 4 + 3x$, hiển thị chúng và tìm nghiệm theo công thức:
\begin{lstlisting}[language=Python]
X = np.random.rand(1000, 1)
y = 4 + 3 * X + .2*np.random.randn(1000, 1) # noise added

# Building Xbar
one = np.ones((X.shape[0],1))
Xbar = np.concatenate((one, X), axis = 1)

A = np.dot(Xbar.T, Xbar)
b = np.dot(Xbar.T, y)
w_lr = np.dot(np.linalg.pinv(A), b)
print('Solution found by formula: w = ',w_lr.T)

# Display result
w = w_lr
w_0 = w[0][0]
w_1 = w[1][0]
x0 = np.linspace(0, 1, 2, endpoint=True)
y0 = w_0 + w_1*x0

# Draw the fitting line
plt.plot(X.T, y.T, 'b.') # data
plt.plot(x0, y0, 'y', linewidth = 2) # the fitting line
plt.axis([0, 1, 0, 10])
plt.show()
\end{lstlisting}

Kết quả:
\begin{lstlisting}
Solution found by formula: w = [[ 4.00305242 2.99862665]]
\end{lstlisting}

Đường thẳng tìm được là đường có màu vàng có phương trình $y \approx 4 + 2.998x$.

\textbf{Tiếp theo} ta viết đạo hàm và hàm mất mát:
\begin{lstlisting}[language=Python]
def grad(w):
 N = Xbar.shape[0]
 return 1/N * Xbar.T.dot(Xbar.dot(w) - y)

def cost(w):
 N = Xbar.shape[0]
 return .5/N*np.linalg.norm(y - Xbar.dot(w), 2)**2;
\end{lstlisting}
\end{vidu}

\paragraph{Kiểm tra đạo hàm.}

\begin{dinhnghia}[title={Numerical gradient}]
Việc tính đạo hàm của hàm nhiều biến thông thường khá phức tạp và rất dễ mắc lỗi, nếu chúng ta tính sai đạo hàm thì thuật toán GD không thể chạy đúng được. Trong thực nghiệm, có một cách để kiểm tra liệu đạo hàm tính được có chính xác không. Cách này dựa trên định nghĩa của đạo hàm (cho hàm một biến):
\[
f'(x) = \lim_{\varepsilon \rightarrow 0}\frac{f(x + \varepsilon) - f(x)}{\varepsilon}
\]
Một cách thường được sử dụng là lấy một giá trị $\varepsilon$ rất nhỏ, ví dụ $10^{-6}$, và sử dụng công thức xấp xỉ hai phía. Cách tính này được gọi là \textbf{numerical gradient}.
\end{dinhnghia}

\begin{congthuc}[title={Xấp xỉ đạo hàm hai phía}]
\[
f'(x) \approx \frac{f(x + \varepsilon) - f(x - \varepsilon)}{2\varepsilon} \qquad (2)
\]
\end{congthuc}

\begin{luuy}[title={Tại sao xấp xỉ hai phía?}]
\textbf{Câu hỏi:} Tại sao công thức xấp xỉ hai phía trên đây lại được sử dụng rộng rãi, sao không sử dụng công thức xấp xỉ đạo hàm bên phải hoặc bên trái?

Có hai cách giải thích cho vấn đề này, một bằng hình học, một bằng giải tích.
\end{luuy}

\paragraph{Giải thích bằng hình học.}

\begin{ynghia}[title={Ba vector xấp xỉ so với đạo hàm chính xác}]
Quan sát hình dưới đây:

\mlcbimg{01_gradient_descent_part1/img08.png}

Trong hình, vector màu đỏ là đạo hàm chính xác của hàm số tại điểm có hoành độ bằng $x_0$. Vector màu xanh lam (có vẻ là hơi tím sau khi convert từ \texttt{.pdf} sang \texttt{.png}) thể hiện cách xấp xỉ đạo hàm phía phải. Vector màu xanh lục thể hiện cách xấp xỉ đạo hàm phía trái. Vector màu nâu thể hiện cách xấp xỉ đạo hàm hai phía. Trong ba vector xấp xỉ đó, vector xấp xỉ hai phía màu nâu là gần với vector đỏ nhất nếu xét theo hướng.

Sự khác biệt giữa các cách xấp xỉ còn lớn hơn nữa nếu tại điểm $x$, hàm số bị bẻ cong mạnh hơn. Khi đó, xấp xỉ trái và phải sẽ khác nhau rất nhiều. Xấp xỉ hai bên sẽ ổn định hơn.
\end{ynghia}

\paragraph{Giải thích bằng giải tích.}

\begin{phuongphap}[title={Khai triển Taylor và sai số $O(\varepsilon)$ vs.\ $O(\varepsilon^2)$}]
Chúng ta cùng quay lại một chút với Giải tích I năm thứ nhất đại học: Khai triển Taylor.

Với $\varepsilon$ rất nhỏ, ta có hai xấp xỉ sau:
\[
f(x + \varepsilon) \approx f(x) + f'(x)\varepsilon + \frac{f''(x)}{2} \varepsilon^2 + \cdots
\]
và:
\[
f(x - \varepsilon) \approx f(x) - f'(x)\varepsilon + \frac{f''(x)}{2} \varepsilon^2 - \cdots
\]
Từ đó ta có:
\[
\frac{f(x + \varepsilon) - f(x)}{\varepsilon} \approx f'(x) + \frac{f''(x)}{2}\varepsilon + \cdots = f'(x) + O(\varepsilon) \qquad (3)
\]
\[
\frac{f(x + \varepsilon) - f(x - \varepsilon)}{2\varepsilon} \approx f'(x) + \frac{f^{(3)}(x)}{6}\varepsilon^2 + \cdots = f'(x) + O(\varepsilon^2) \qquad (4)
\]
Từ đó, nếu xấp xỉ đạo hàm bằng công thức $(3)$ (xấp xỉ đạo hàm phải), sai số sẽ là $O(\varepsilon)$. Trong khi đó, nếu xấp xỉ đạo hàm bằng công thức $(4)$ (xấp xỉ đạo hàm hai phía), sai số sẽ là $O(\varepsilon^2) \ll O(\varepsilon)$ nếu $\varepsilon$ nhỏ.

Cả hai cách giải thích trên đây đều cho chúng ta thấy rằng, xấp xỉ đạo hàm hai phía là xấp xỉ tốt hơn.
\end{phuongphap}

\paragraph{Với hàm nhiều biến.}

\begin{luuy}[title={Áp dụng numerical gradient cho nhiều biến}]
Với hàm nhiều biến, công thức $(2)$ được áp dụng cho từng biến khi các biến khác cố định. Cách tính này thường cho giá trị khá chính xác. Tuy nhiên, cách này không được sử dụng để tính đạo hàm vì độ phức tạp quá cao so với cách tính trực tiếp. Khi so sánh đạo hàm này với đạo hàm chính xác tính theo công thức, người ta thường giảm số chiều dữ liệu và giảm số điểm dữ liệu để thuận tiện cho tính toán. Một khi đạo hàm tính được rất gần với numerical gradient, chúng ta có thể tự tin rằng đạo hàm tính được là chính xác.
\end{luuy}

\begin{vidu}[title={Code kiểm tra đạo hàm}]
Dưới đây là một đoạn code đơn giản để kiểm tra đạo hàm và có thể áp dụng với một hàm số (của một vector) bất kỳ với \texttt{cost} và \texttt{grad} đã tính ở phía trên.
\begin{lstlisting}[language=Python]
def numerical_grad(w, cost):
 eps = 1e-4
 g = np.zeros_like(w)
 for i in range(len(w)):
 w_p = w.copy()
 w_n = w.copy()
 w_p[i] += eps
 w_n[i] -= eps
 g[i] = (cost(w_p) - cost(w_n))/(2*eps)
 return g

def check_grad(w, cost, grad):
 w = np.random.rand(w.shape[0], w.shape[1])
 grad1 = grad(w)
 grad2 = numerical_grad(w, cost)
 return True if np.linalg.norm(grad1 - grad2) < 1e-6 else False

print( 'Checking gradient...', check_grad(np.random.rand(2, 1), cost, grad))
\end{lstlisting}

Kết quả:
\begin{lstlisting}
Checking gradient... True
\end{lstlisting}

(Với các hàm số khác, người học chỉ cần viết lại hàm \texttt{grad} và \texttt{cost} ở phần trên rồi áp dụng đoạn code này để kiểm tra đạo hàm. Nếu hàm số là hàm của một ma trận thì chúng ta thay đổi một chút trong hàm \texttt{numerical\_grad}, có thể chỉnh sửa tương tự).

Với bài toán Linear Regression, cách tính đạo hàm như trong $(1)$ phía trên được coi là đúng vì sai số giữa hai cách tính là rất nhỏ (nhỏ hơn $10^{-6}$). Sau khi có được đạo hàm chính xác, chúng ta viết hàm cho GD:
\begin{lstlisting}[language=Python]
def myGD(w_init, grad, eta):
 w = [w_init]
 for it in range(100):
 w_new = w[-1] - eta*grad(w[-1])
 if np.linalg.norm(grad(w_new))/len(w_new) < 1e-3:
 break
 w.append(w_new)
 return (w, it)

w_init = np.array([[2], [1]])
(w1, it1) = myGD(w_init, grad, 1)
print('Solution found by GD: w = ', w1[-1].T, ',\nafter %d iterations.' %(it1+1))
\end{lstlisting}

Kết quả:
\begin{lstlisting}
Solution found by GD: w = [[ 4.01780793 2.97133693]] ,
after 49 iterations.
\end{lstlisting}

Sau $49$ vòng lặp, thuật toán đã hội tụ với một nghiệm khá gần với nghiệm tìm được theo công thức.

Dưới đây là hình động minh họa thuật toán GD.

\mlcbimg{01_gradient_descent_part1/img09.png}
\mlcbimg{01_gradient_descent_part1/img10.png}
\end{vidu}

\begin{ynghia}[title={Hai panel minh họa GD cho Linear Regression}]
Trong hình bên trái, các đường thẳng màu đỏ là nghiệm tìm được sau mỗi vòng lặp.

Trong hình bên phải, đây là khái niệm: \textbf{đường đồng mức}.
\end{ynghia}

\subsubsection{Đường đồng mức (level sets)}

\begin{dinhnghia}[title={Level sets}]
Với đồ thị của một hàm số với hai biến đầu vào cần được vẽ trong không gian ba chiều, nhiều khi chúng ta khó nhìn được nghiệm có khoảng tọa độ bao nhiêu. Trong toán tối ưu, người ta thường dùng một cách vẽ sử dụng khái niệm \textbf{đường đồng mức} (level sets).

Trong bản đồ tự nhiên trong các bản đồ tự nhiên, để miêu tả độ cao của các dãy núi, người ta dùng nhiều đường cong kín bao quanh nhau như sau:

Ví dụ về đường đồng mức trong các bản đồ tự nhiên. (Địa lý 6: Đường đồng mức là những đường như thế nào?)

Các vòng nhỏ màu đỏ hơn thể hiện các điểm ở trên cao hơn.

Trong toán tối ưu, người ta cũng dùng phương pháp này để thể hiện các bề mặt trong không gian hai chiều.
\end{dinhnghia}

\begin{ynghia}[title={Level sets trong ví dụ Linear Regression}]
Quay trở lại với hình minh họa thuật toán GD cho bài toán Linear Regression bên trên, hình bên phải là hình biểu diễn các level sets (đường đồng mức). Tức là tại các điểm trên cùng một vòng, hàm mất mát có giá trị như nhau. Trong ví dụ này, hình minh họa giá trị của hàm số tại một số vòng. Các vòng màu xanh có giá trị thấp, các vòng tròn màu đỏ phía ngoài có giá trị cao hơn. Điểm này khác một chút so với đường đồng mức trong tự nhiên là các vòng bên trong thường thể hiện một thung lũng hơn là một đỉnh núi (vì chúng ta đang đi tìm giá trị nhỏ nhất).
\end{ynghia}

\begin{vidu}[title={Learning rate nhỏ hơn trên Linear Regression}]
Thử nghiệm với learning rate nhỏ hơn, kết quả như sau:

\mlcbimg{01_gradient_descent_part1/img12.png}
\end{vidu}

\begin{luuy}[title={Hội tụ chậm khi $\eta$ quá nhỏ}]
Tốc độ hội tụ đã chậm đi nhiều, thậm chí sau $99$ vòng lặp, GD vẫn chưa tới gần được nghiệm tốt nhất. Trong các bài toán thực tế, chúng ta cần nhiều vòng lặp hơn $99$ rất nhiều, vì số chiều và số điểm dữ liệu thường là rất lớn.
\end{luuy}

%------------------------------------------------------------------------------
\subsection{1.4. Ví dụ: hai cực tiểu phụ thuộc khởi tạo}
%------------------------------------------------------------------------------

\begin{vidu}[title={Hàm $f(x,y) = (x^2 + y - 7)^2 + (x - y + 1)^2$}]
Trước khi chuyển sang Chương~2, xét thêm một ví dụ hàm hai biến:

Hàm số $f(x, y) = (x^2 + y - 7)^2 + (x - y + 1)^2$ có hai điểm local minimum màu xanh lục tại $(2, 3)$ và $(-3, -2)$, và chúng cũng là hai điểm global minimum. Trong ví dụ này, tùy vào điểm khởi tạo mà chúng ta thu được các nghiệm cuối cùng khác nhau.

\mlcbimg{01_gradient_descent_part1/img14.png}
\end{vidu}

%------------------------------------------------------------------------------
\subsection{1.5. Kết chương 1}
%------------------------------------------------------------------------------

\begin{tongket}[title={Tóm tắt Chương 1}]
Chương này xây dựng \textbf{Gradient Descent} từ quan sát đạo hàm một biến đến cập nhật vector $\boldsymbol{\theta}$, với minh họa Python và Linear Regression. Các ý then chốt: đi \textbf{ngược gradient}; $\eta$ và $\boldsymbol{\theta}_0$ quyết định hội tụ; kiểm tra $\nabla \mathcal{L}$ bằng numerical gradient; đường đồng mức giải thích đường đi zigzag. Chương~2 bổ sung cách scale dữ liệu lớn và tăng tốc hội tụ.
\end{tongket}

%------------------------------------------------------------------------------
\subsection{Tài liệu tham khảo}
%------------------------------------------------------------------------------

\begin{phuongphap}[title={Tài liệu đọc thêm}]
\begin{enumerate}
 \item Sebastian Ruder, \emph{An overview of gradient descent optimization algorithms} --- \url{https://ruder.io/optimizing-gradient-descent/}
 \item \emph{An Interactive Tutorial on Numerical Optimization} --- \url{https://www.benfrederickson.com/numerical-optimization/}
 \item Andrew Ng, \emph{Gradient Descent} (Coursera Machine Learning) --- \url{https://www.coursera.org/learn/machine-learning}
\end{enumerate}
\end{phuongphap}
